\documentclass{article}
\usepackage[T1]{fontenc}
\usepackage{lmodern}
\usepackage{graphicx}
\usepackage{amsmath,amssymb}
\usepackage[a4paper,margin=2cm]{geometry}
\usepackage[absolute,overlay]{textpos}
\setlength{\TPHorizModule}{1mm}
\setlength{\TPVertModule}{1mm}
\usepackage[indonesian]{babel}
\usepackage{hyperref}
\setlength{\emergencystretch}{2em}
% unit: Halaman 12

\begin{document}

% === Halaman 12 (python/native) ===

Jenis-jenis steganalisis

1. Targeted steganalysis

• Teknik steganalisis yang bekerja pada algoritma steganografi spesifik, dan kadang-

kadang dibatasi hanya pada format media tertentu saja.

• Teknik ini mempelajari dan menganalisis algoritma penyisipan, lalu menemukan 

statistik yang berubah setelah penyisipan. 

• Hasil steganalisis sangat akurat, tetapi tidak fleksibel karena tidak dapat diperluas 

untuk algoritma steganografi yang lain atau format media yang berbeda.

12

\end{document}
